%% The first command in your LaTeX source must be the \documentclass command.
%%
%% Options:
%% twocolumn : Two column layout.
%% hf: enable header and footer.
\documentclass[
  % twocolumn,
  % hf,
]{ceurart}

%%
%% One can fix some overfulls
\sloppy


\usepackage{caption}

%%
%% Minted listings support
%% Need pygment <http://pygments.org/> <http://pypi.python.org/pypi/Pygments>
\usepackage{listings}
%% auto break lines
\lstset{breaklines=true}

\usepackage{hyperref}
\usepackage{svg}

%%
%% end of the preamble, start of the body of the document source.
\begin{document}

\graphicspath{{./graphics}}

%%
%% Rights management information.
%% CC-BY is default license.
\copyrightyear{2025}
\copyrightclause{Copyright for this paper by its authors.
  Use permitted under Creative Commons License Attribution 4.0
International (CC BY 4.0).}


\conference{Workshop on Computational Drama Analysis, 3 September
2025, Berlin, DE}

\title{The Measure of a Messenger: Meter as potential avenue of
 differentiation}

%%
%% The "author" command and its associated commands are used to define
%% the authors and their affiliations.
\author[1]{Charles Pletcher}[%
  orcid=0000-0003-2734-5511,
  email=charles.pletcher@tufts.edu,
  url=https://github.com/pletcher,
]
\cormark[1]
\address[1]{Tufts University, Department of Classical Studies, 10
Upper Campus Road, Medford, Massachusetts 02155, United States}

%% Footnotes
\cortext[1]{Corresponding author.}

%%
%% The abstract is a short summary of the work to be presented in the
%% article.
\begin{abstract} This paper searches for a metrical definition of
 messenger-type speech by way of arguing that such speeches (\textit
 {Botenberichte}) in the ancient tragedies of Aeschylus, Sophocles, and
 Euripides have not received the computational attention that they deserve,
 but also that such attention is impossible given the current dearth of data
 for Ancient Greek tragedy. I first draw attention to this gap by surveying
 traditional literary approaches to the genre, showing how the definition of
 `messenger speech' is already problematic from a qualitative perspective.
 Next, I present two quantitative approaches to the meter of messenger speeches,
 neither of which manages to identify any distinguishing characteristics. In
 conclusion, I offer a brief discussion of these preliminary findings and
 potential avenues for further inquiry.
\end{abstract}

%%
%% Keywords. The author(s) should pick words that accurately describe
%% the work being presented. Separate the keywords with commas.
\begin{keywords}
  tragedy \sep
  messenger speeches \sep
  iambic trimeter \sep
  metrical resolution
\end{keywords}

%%
%% This command processes the author and affiliation and title
%% information and builds the first part of the formatted document.
\maketitle

\section{Introduction}

This paper asks whether we can detect a `metrical signature' for messengers'
speeches\footnote{Throughout this paper, I use `messenger-type speech' to refer to
the category of speeches under investigation. But `messenger' describes a heterogenous
class of individuals who do not necessarily conform to a single role or character type.
I thus also use the term `news-bringer' when I mean to include characters who deliver
news but who are not generally considered `messengers.' Further, the analyses in this paper
} in ancient Athenian tragedy. It pursues this line of inquiry through
quantitative analyses of the meter in the three canonical tragedians
(Aeschylus, Sophocles, and Euripides). In so doing, this paper sheds new
light on questions that have previously been addressed mainly from
qualitative angles, challenging the prevailing view that we can easily
distinguish messenger-type speech from other iambic trimeters (the main
spoken meter of tragedy).

In the first pages of \textit{Narrative in Drama: The Art of the Euripidean
Messenger-Speech}, Irene J. F. de Jong offers the following definition of
(Euripidean) messenger-speeches: `The Euripidean messenger-speech or ἀγγελικὴ
ῥῆσις is a long, continuous speech in which a messenger reports events that
have taken place off-stage. The content of this speech is a narrative, more
particularly a first-person narrative…' \cite[1–2]{deJong1991}. She goes on to
argue that Euripides makes signficantly greater use of first-person narrative
than Sophocles and Aeschylus, citing \cite[2 n. 4]{deJong1991}

\begin{itemize}
  \item Aeshylus, \textit{Persians} 480–514
  \item Aeschylus, \textit{Agamemnon} 636–80
  \item Sophocles, \textit{Antigone} 407–40, 1192–1243
  \item Sophocles, \textit{Oedipus Tyrannus} 1237–85
  \item Sophocles, \textit{Oedipus Colonus} 1586–1666
\end{itemize}

As the only examples of first-person narrative outside of Euripides' oeuvre.
Although de Jong takes great pains (and around sixty pages) establishing the
first-person focalization of Euripides' messenger-speeches, a cursory glance
at, for example, the Tutor's speech in Sophocles' \textit{Electra} (680–763)
casts enough doubt on de Jong's criteria that it raises further questions about
not only Euripidean messenger-speech exceptionalism but messenger-speech as a
category in its own right. \footnote{The Tutor's speech is rightly considered
  an exemplary `messenger-type speech' in its own right (so
  \cite{Marshall2006} and \cite{Kells1973}, among others), and it is largely
  for this reason that its apparent violation of Euripidean messenger-speech
exceptionalism raises other questions.} In this paper, I present two
computational null results that cast further doubt on messengers'
exceptionalism. By showing how ordinary messengers' speeches are, I mean both
to highlight their dramatic---as opposed to narrative---place and to
demonstrate how their performances contribute to tragedy's dramatic
negotiations for truth.\footnote{I use `messenger-speech' when referring to
  de Jong's concept, and `messenger's speech,' `messengers' speeches,' or
  `messenger-type speech' when referring to the general category of dramatic
monologue.}

The paper begins with a brief and necessarily incomplete overview of the
massive (and growing) bibliography on messengers and their speeches. I then
discuss the methodologies and experiments by which this paper's results were
obtained, beginning with the library that was used for scanning iambic
trimeters, the spoken meter of Greek tragedy and thus the meter for most
messenger-type (Euripides singing Phrygian in \textit{Orestes} marking a
notable exception). The first two experiments examine two interrelated ways in
which meter might be supposed to distinguish messengers' speeches from other
characters' speech, but which instead show that meter does not sufficiently
distinguish messengers from other tragic speakers. The third experiment uses
TF--IDF and 4-character-grams to cluster the approximately 17,000 lines of
iambic trimeter in tragedy with all of the speeches from the Homeric corpus,
obtained through the Digital Initiative for Classics: Epic Speeches (DICES)
API. These experiments and their results are far from exhaustive, but they
offer a new way of understanding the place of messengers in drama---a place,
moreover, that appears no less `dramatic' than their traditionally less
`narrative' peers.

\section{A brief history of \textit{Botenberichte}}

The usual history of (so-called) `messenger speeches' in drama begins with
Homer: the ancient bard's narrative \textit{enargeia}, to borrow a term from
Pseudo-Longinus, is said to animate the off-stage events that news-bringers
narrate. De Jong has perhaps most famously applied a narratological lens to
both Homer and tragic messengers, followed a few years later by Barbara
Goward and James Barrett with their own seminal works on messenger speeches
in 2002 \cite{deJong2011, deJong1991, Goward1999, Barrett2002}. De Jong
championed the view that messengers are not featureless conduits for the
dramatist's voice but are instead focalizers with bias and selective memory.
Barrett then qualified this view by noting that messengers' narrative
effectiveness is inversely proportional to their stage presence: `The
messenger's own body,' he writes, `is implicitly—and explicitly, on
occasion—a hindrance, inasmuch as his success as a messenger depends upon his
acquiring the invulnerability granted by freedom from bodily
constraints' \cite[100]{Barrett2002}. More recently, Margaret Dickin showed
how the roles of messengers in tragedy increased over time and became
`vehicles for star-power,' while others, like Felix Budelmann and Evert van
Emde Boas, explored how messengers direct the audience's affective
energies \cite{Dickin2009, Budelmann.vanEmdeBoas2020}. Florence Yoon's
interventions have further disrupted the `category' of messengers by
revealing the mistakes of painting with too broad a brush:
\textit{angeloi}, messengers proper, and \textit{kerykes}, `heralds,'
perform different functions in tragedy \cite{Yoon2022}.

The chief difficulty running through much of the literature on messenger
speeches has to do with defining who, precisely, is a news-bringer\footnote
{Along with ἄγγελοι, `messengers,' I also count κήρυκες, `heralds,' and
various other unnamed deliverers of information. For the purposes of the
computational studies that follow, however, I divide these speakers into
different classes, described below. As I hope to show, any metrical
distinctions between classes of characters in tragedy are artificial at
best.} and what, precisely, is a messenger speech. Part of the difficulty
arises from the fact that named characters at times appear to deliver
speeches that resemble messenger speeches: nearly every named character in
Sophocles'
\textit{Trachiniae} does so, as Sir Richard Jebb noted over a century ago
\cite{Jebb1902}; and even Clytaemestra's famous "Beacon Speech" in Aeschylus'
\textit{Agamemnon} resembles speeches of news-bringers in its
subject-matter. Thus, some scholars insist on anonymity as a prime criterion
for the deliverer of a messenger speech, but such an insistance runs the risk
of obscuring the similarities that so-called `messenger speeches' have with
respect to other monologues in drama. These similarities matter, I argue,
because they highlight the various strategies that drama uses to interrogate
the fraught and problematic delivery of information that the recipients cannot
verify. The tragedians, to state the matter bluntly, offer some of the earliest
extended encounters with `fake news.'

Yet despite the steady stream of innovative approaches to messenger speeches
in tragedy, and despite Simon Perris' call for a `usable, robust, formal
working definition' of messenger speech, no clear qualitative definition has
emerged.\cite{Perris2011} On the quantitative side, there have been few
attempts to address the various classes of news-bringers from a computational
angle. Where such examinations have occurred, they have, to my knowledge,
taken place in the context of broader computational studies of drama, such as
the work of Julia Jennifer Beine, Frank Fischer, and Viktor Illmer presented
at the Digital Humanities Conference in 2024; the examination of the chorus
by María Teresa Santa María Fernández and Monika Dabrowska; or the early work
on social network theory and Greek tragedy by Jeff Rydberg-Cox.\cite
{Beine.etal2024, Fernandez.Dabrowska2023, Rydberg-Cox2011}

\section{Messengers and Metre}

With the notable exception of the singing Phrygian in Euripides' \textit
{Orestes}, most news-bringers in tragedy speak in iambic trimeters. The most
basic form of iambic trimeter, x — u — | x — u — | x — u —, lends itself to a
striking amount of variation: nearly every long syllable (—) can have two
short syllables in its place, provided several additional rules are heeded.
Such substitutions are called `resolutions,' as the long syllable is
understood to `resolve' into two short syllables. The vertical lines in the
schema above separate the dipodes or dipodies, literally the `two-foot`
groupings of iambs. At \textit{Poetics} 1449 a 24, Aristotle dubs iambic `the
most speech-like of the metres,' no doubt due to its flexibility.

Recent studies on the meters of tragedy have either focused on broad
methodological challenges, such as the recent thesis by Erik Henrikkson; or
they have dealt with more directed questions like L. P. E. Parker's study of
the impact of resolution-avoidance on the vocabulary of Sophocles'
trimeters.\cite{Henriksson, Parker2022} A recent study by Boris Maslov also
examines the interplay of rhythm and vocabulary in trimeters, using the
Prosimetron database and toolkit to analyze Sophocles' \textit
{Ajax} and \textit{Oedipus Tyrannus} alongside the Byzantine \textit
{Christus Patiens}, a `Christian tragedy' modelled on Euripides' \textit
{Bacchae}.\cite{Maslov2022} In a similar vein, Felipe G. Hernández Muñoz's
recent paper looks at the placement of the pronoun ἐγώ in iambic trimeter and
its potential relationship with speaker class (royalty, slave, etc.).\cite
{HernandezMunoz2019} To my knowledge, although resolution in iambic trimeter
has been studied extensively from a descriptive point of view, and has
notably been used for dating the plays of Euripides, no inquiry has been made
into the possibility that it might be used differently in different classes
of speakers.\footnote{See \cite{Parker2022} p. 232 n. 3 for additional
bibliography on iambic trimeters in tragedy.}

In this paper, I undertake two quantitative approaches to meter in tragedy,
with a focus on whether messengers and their ilk exhibit a discernible
metrical signature by which they might be differentiated from other speakers.
I investigate two specific phenomena: (1) \textbf{the number of resolutions
per line of each character's speech}, that is, the number of times that a
single long syllable in an iamb is replaced by two short syllables; and
(2) \textbf{the placement of resolutions in characters' trimeters}, that is,
whether resolution appears in position(s) 2, 4, 6, 8, and/or 10 in a given
line.\footnote{Caesurae were also analyzed as part of these experiments, but
as with the other experiments, they showed no substantial correlation with
speaker class.}

At the risk of discouraging the reader from continuing, both experiments turn
up null or nearly null results, with extremely weak correlation at best: the
analyzed features are insufficient to distinguish messengers from other
classes of tragic characters. Yet the reasons for this lack of differentation
can still lead to a more nuanced understanding of the \textit{dramatic} place
of messengers in tragedy, as opposed to their relegation to `\textit
{narrative} in drama.' By exploring the ways that we can measure a
character's `messenger-type-ness,' or the extent to which they perform a
`messenger function' (to borrow a term from C. W. Marshall \cite
{Marshall2006}), we shed new light on this often-misunderstood class of
dramatic actors and thus better understand how they reflect their culture's
relationship to secondhand information.

\section{Data preparation}

Data were prepared using TEI XML files from DraCor and the Perseus Project.
For determining scansion, a rules-based library was built in the Elixir
programming language. Although the deviation from Python might seem unusual,
Elixir has the benefit of platform stability---anecdotally, Python projects
break approximately every six months, while I have been able to maintain
Elixir-based experiments for nearly a decade without major interruption---as
well as first-class support for multi-threaded and parallel computing,
techniques that will be beneficial for larger corpora. This early
programming-language decision will mean straightforward parallelization as/if
the corpus grows. It also means that these tools are reusable for larger
projects, even if the developers of those projects are not familiar with
Elixir: they need only install the Elixir runtime and run the pre-written
processing scripts to output their own scansion analysis.

Experiments were run in Livebook, an Elixir-based corollary to Jupyter
notebooks. The relevant Livebooks and supporting files are available upon
request. (They are currently hosted on a platform other than GitHub and so
cannot be anonymized.)

For the sake of efficiency, the rules-based scansion parser was initially
built with the assistance of Claude Code. While some of that code remains in
the final implementation, the author found several major bugs that were
leading to under-counting the number of trimeters in characters'
speech---trimeter should make up about 60\% of tragic dialogue, but the
`vibe-coded' implementation of the metrical analyser only captured 40\%.
Thus, the author rewrote much of the code to fix these errors, and the
library now reports valid trimeters as 57\% of lines scanned. The author
verified the scansion results by examining key passages from each of the
three playwrights, and the author wrote all of the experiments and conducted
the analysis without AI assistance.

The scansion parser first syllabifies a text and then parses line by line to
determine whether the line can scan as valid iambic trimeter. Out of 44,482
lines of tragedy in the corpus, 26,433 scan as valid trimeters; 81 were
skipped due to parsing errors. The corpus of iambic trimeters thus represents
about 57.17\% of the lines of tragedy in the corpus, matching the rough
division between odes and episodes in most plays, which hovers around 60\%.

For all of the experiments, each speaker was labeled with one of eight
`speaker classes':(1) chorus, (2) divine, (3) herald, (4) messenger,
(5) other, (6) prophet,(7) royalty, and (8) slave (in alphabetical order).
These labels were first applied automatically to unambiguous characters---all
Ἄγγελοι were labeled `messenger,' for example. Subsequently, a list of
ambiguous characters was generated and labeled by hand.

The definition of messenger for the purposes of labeling was deliberately
strict. Although, as noted above, the category of messenger-type speech has
notoriously fuzzy boundaries, the labeling for these experiments followed
character naming conventions exclusively. We should note that this labeling
method is still artificial: while we cannot know specifics about when line
assignemnt labels entered the manuscript tradition, our earliest manuscripts
include only paragraphoi to signal a change of speaker, but they do not label
the speakers specifically.\footnote{See \cite
{Mastronarde2016}.} Nevertheless, it seemed best to use what are now the
canonical labels for speaker class, rather than to rely on necessarily
subjective judgment while attempting to establish quantitative parameters for
identifying messenger-type speech.

\section{Parsing iambic trimeter}

The rules-based parser works as follows (visualized also in the accompanying
above): (1) First, the raw Greek text is read to determine whether it is Beta
Code or polytonic Greek. These experiments work exclusive with polytonic
Greek, so this step is essentially a noop. (2) Second, the raw text is
tokenized by grapheme (letters with diacritical markings attached).
(3) Subsequently, the graphemes are grouped by syllables, accounting for
diphthongs, elision, and double consonants. Compounds such as εἰσ-φέρω are
not explicitly handled, as doing so would add out-of-scope morphological
analysis to this scanner with minimal improvement on the scansion analysis:
for example, in the case of εἰσ-φέρω, εἰ- will always scan long, and the
epsilon in φέρω will always scan short, whether the compound is grouped
correctly as εἰσ-φέρω or incorrectly as εἰ-σφέρω. (4) Syllables are
classified as short, long, anceps, or unknown according to standard rules of
Ancient Greek scansion. According to the usual convention in tragic
verse, \textit{muta cum liquida} (a stop consonant followed by a lambda or a
rho) does not result in the usual lengthening that applies to other short
vowels followed by two consonants. Other double-consonant clusters will cause
preceding short vowels to be scanned long by this scanner.\footnote
{Unknown vowel quantities are logged in a dictionary for later manual
annotation. This annotation process has not yet been completed.} (5) The
scanned syllables are mapped onto the iambic trimeter schema of three
dipodies. (6) Features such as resolution and dactylic events are labeled
according to the flow chart below.

\noindent
\begin{minipage}{\linewidth}
\makebox[\linewidth]{
\begin{tabular}{c}
(1) raw text \\
\downarrow{} \\
(2) grapheme tokenization \\
\downarrow{} \\
(3) syllabification \\
\downarrow{} \\
(4) quantification (labeling of short, long, and anceps syllables) \\
\downarrow{} \\
(5) parsing (mapping to iambic trimeter) \\
\downarrow{} \\
(6) feature extraction (label resolutions and dactylic events)
\end{tabular}}
\end{minipage}

\noindent
Dactylic events refer specifically to the possibility of a long anceps
followed by a resolved iambic `up-beat,' i.e., the familiar epic — u u
(long short short). A hypothesis concerning the `epic-ness' of messengers
in tragedy might suppose that they make greater use of dactyls than other
classes of speakers. This hypothesis turns out false, however.


\section{Resolution count}

We begin by counting resolutions in iambic trimeters. The frequency of
resolution offers a sense of how often each character class adds more
syllables than are strictly necessary to a line of iambic trimeter. The
strength of the relationship between speaker class and resolution count was
tested using [Harald] Cramér's $φ_c$ (also known as Cramér's V). As Cramér
writes, `The correlation coefficient and the correlation ratio both server
to characterize, in the sense explained above, the `degree of dependence'
between two variables.... $φ^2 = 0$ when and only when the variables are
independent.' \cite[282]{Cramer2016} Cramér goes on to explain that $φ_c$
will range between \(0\) (no correlation) and \(1\) (full corrleation). Thus,
by measuring Cramér's $φ_c$ for speaker class and number of resolutions, we
can determine whether the two are related and the strength of that relation.

For 8 speaker classes, the correlation with resolution count was
0.0437---meaning a miniscule connection between the two variables. The figure
below shows the relative frequency of resolutions per total lines by speaker
class.

\begin{center}
\begin{tabular}{c | c | c | c}
\textbf{speaker class} & \textbf{resolutions} & \textbf{total lines} & \textbf{rel. freq. resolved} \\
prophet & 89 & 350 & 0.2543 \\
other & 270 & 1478 & 0.1827 \\
chorus & 289 & 1685 & 0.1715 \\
messenger & 383 & 2336 & 0.1640 \\
royalty & 2445 & 16205 & 0.1509 \\
divine & 265 & 1927 & 0.1375 \\
slave & 83 & 893 & 0.0929 \\
herald & 1 & 665 & 0.0859
\end{tabular}
\end{center}

\section{Resolution placement}

As with total resolutions, resolution placement has a negligible positive
correlation with speaker class, peaking in positions 6 (0.0924) and 2 (0.0622).
These correlations are meaningful, but they are too small with our limited data
to suggest predictive value.


\begin{center}
\begin{tabular}{c | c}
\textbf{position} & \textbf{Cramér's $φ_c$} \\
2 & 0.0622 \\
4 & 0.0468 \\
6 & 0.0924 \\
8 & 0.0583 \\
10 & 0.0184
\end{tabular}
\end{center}


Indeed, we can see small but meaningful differences in the absolute and
relative frequencies for each resolution for every position but 10. (See full
results at >>REDACTED FOR ANONYMITY<<) These differences, however, hardly provide any
descriptive value. The relative frequencies of resolution are so low for each
position---generally under 10\%--- that it is difficult if not impossible to
infer a metrical signature for any class of speaker from data drawn from the
extant plays.

Below are visualized the relative frequencies of resolution by position and by
tragedian. The breakdown by tragedian is necessary because of the state of
the extant corpus: 19 (counting, perhaps contraversially, \textit{Rhesus}) of
Euripides' plays survive, compared to 7 each for Aeschylus and Sophocles,
plus fragments from Sophocles' satyr play \textit{Trackers}. Although, as
expected, Aeschylus and Sophocles use resolution less frequently than
Euripides (relative to total lines written), the unbalanced representation in
the corpuse leads to greater observed variability in the distribution of
resolutions among the speaker classes---meaning, for example, that divinities
appear to lead the relative frequency of resolutions in the sixth(5.7\%) and
eighth (4.3\%) positions, yet there are only two speaking divinities present
in Sophoclean trimeter: Athena in the prologue of \textit{Ajax} and
Herakles \textit{ex machina} in \textit{Philoctetes}. Relative frequencies
are visualized both as a proportion of each speaker class's resolution
placement and as a grouped bar chart across speaker classes.


\begin{center}
\captionof{figure}{Graphic of relative frequencies by resolution placement for Aeschylus}\label{fig:aeschylus-rel-freq}
\includesvg{aeschylus-rel-freq}
\end{center}

\begin{center}
\captionof{table}{Table of relative frequencies by resolution placement for Aeschylus}
\begin{tabular}{c | c | c | c | c | c} \\
\textbf{speaker class} & resolves at 02 & resolves at 04 & resolves at 06 & resolves at 08 & resolves at 10 \\
chorus & 0.0197 & 0.0108 & 0.0394 & 0.0269 & 0.0161 \\
divine & 0.0043 & 0.0021 & 0.0361 & 0.0096 & 0.0021 \\
herald & 0.0063 & 0.0000 & 0.0438 & 0.0188 & 0.0000 \\
messenger & 0.0081 & 0.0081 & 0.0569 & 0.0352 & 0.0081 \\
other & 0.0160 & 0.0000 & 0.0080 & 0.0000 & 0.0080 \\
royalty & 0.0120 & 0.0036 & 0.0455 & 0.0150 & 0.0006 \\
slave & 0.0000 & 0.0000 & 0.0488 & 0.0000 & 0.0000 \\
\textbf{std. dev.} & 0.0069 & 0.0043 & 0.0155 & 0.0132 & 0.0061 \\
\end{tabular}
\end{center}

\begin{center}
\captionof{figure}{Graphic of frequencies by resolution placement for Sophocles}\label{fig:sophocles-rel-freq}
\includesvg{sophocles-rel-freq}
\end{center}

\begin{center}
\captionof{table}{Table of relative frequencies by resolution placement for Sophocles}
\begin{tabular}{c | c | c | c | c | c} \\
\textbf{speaker class} & resolves at 02 & resolves at 04 & resolves at 06 & resolves at 08 & resolves at 10 \\
chorus & 0.0447 & 0.0419 & 0.0559 & 0.0279 & 0.0168 \\
divine & 0.0000 & 0.0000 & 0.0571 & 0.0429 & 0.0000 \\
herald & 0.0000 & 0.0000 & 0.0100 & 0.0000 & 0.0100 \\
messenger & 0.0117 & 0.0067 & 0.0417 & 0.0150 & 0.0017 \\
other & 0.0086 & 0.0172 & 0.0086 & 0.0086 & 0.0000 \\
prophet & 0.0244 & 0.0000 & 0.0163 & 0.0163 & 0.0000 \\
royalty & 0.0151 & 0.0075 & 0.0273 & 0.0101 & 0.0019 \\
slave & 0.0090 & 0.0045 & 0.0495 & 0.0090 & 0.0000 \\
\textbf{std. dev.} & 0.0146 & 0.0142 & 0.0203 & 0.0134 & 0.0062 \\
\end{tabular}
\end{center}

\begin{center}
\captionof{figure}{Graphic of relative frequencies by resolution placement for Euripides}\label{fig:euripides-rel-freq}
\includesvg{euripides-rel-freq}
\end{center}

\begin{center}
\captionof{table}{Table of relative frequencies by resolution placement for Euripides}
\begin{tabular}{c | c | c | c | c | c} \\
\textbf{speaker class} & resolves at 02 & resolves at 04 & resolves at 06 & resolves at 08 & resolves at 10 \\
chorus & 0.0455 & 0.0416 & 0.0754 & 0.0247 & 0.0195 \\
divine & 0.0524 & 0.0437 & 0.0983 & 0.0295 & 0.0022 \\
herald & 0.0134 & 0.0134 & 0.0635 & 0.0234 & 0.0033 \\
messenger & 0.0369 & 0.0328 & 0.0990 & 0.0411 & 0.0021 \\
other & 0.0479 & 0.0342 & 0.0889 & 0.0368 & 0.0068 \\
prophet & 0.1013 & 0.0661 & 0.1189 & 0.0573 & 0.0176 \\
royalty & 0.0417 & 0.0335 & 0.0972 & 0.0309 & 0.0046 \\
slave & 0.0190 & 0.0254 & 0.0444 & 0.0143 & 0.0000 \\
\textbf{std. dev.} & 0.0267 & 0.0153 & 0.0235 & 0.0130 & 0.0074 \\
\end{tabular}
\end{center}

As is readily visible in these figures, relative resolution count varies
significantly by dramatist, with Euripides' maximum relative frequency
(11.9\% for prophets) more than doubling the maxima of Sophocles (5.7\% for
divinities) and Aeschylus (5.7\% for messengers). Aeschylus and Sophocles
exhibit more variable resolution placement among speaker classes than
Euripides, but this observation is likely due to the smaller sample size for
these tragedians: the data include Aeschylus' 7 canonical plays
(including \textit{Prometheus Bound}) and Sophocles' 7 plays and the
substantial fragments of the satyr play (\textit{Trackers}); but the data
also include all 17 of Euripides' tragedies and 2 of his satyr plays(\textit
{Cyclops} and \textit{Alcestis}), meaning that we have significantly more
information for his plays than for the plays of the other two tragedians. The
marked differences in distributions between the three tragedians' resolution
placements supports separating our analyses; but at the same time, the
smaller sample sizes for Aeschylus and Sophocles limit our ability to draw
inferences from the resolution distributions.


\section{Discussion}


These experiments show that there \textit{are} detectable metrical differences
among speaker classes in tragedy, perhaps especially among messengers in
Euripides, where resolutions at positions 6 and 8 are relatively frequent
while resolutions in position 10 are almost entirely avoided. Compare this
distribution to the distribution of resolutions in the Euripidean chorus, who
use resolutions much more frequently throughout their trimeters.(Note, again,
that these experiments do not look at the chorus' lyric metres.) However,
these experiments also point up the impossibility of making meaningful
classifications based on the available data: we could not for instance, take
an unassigned line from the tragic fragments and determine whether it was a
messenger's or not. Even in experiments where we limited speaker classes to
`messenger' and `non-messenger,' classification algorithms achieved greater
than 90\% accuracy simply by labeling every line as `non-messenger,' because
messengers \textit{strictu sensu} generally speak fewer than 10\% of the
lines in a a given tragedy. These null quantitative results should not be
understood as contradicting the more literary analyses of de Jong and others.
This paper's results only suggest that meter fails as a distinguishing
feature of messenger-type speech; but other features, such as lexis and
syntax, might still reveal key differences between messengers and other types
of characters.

These differences matter because messengers model how audiences understand,
interpret, and evaulate second-hand news and the people who bring it to them.
But in the case of this paper, the \textit{lack} of apparent metrical
differences matters, too. Messengers, according to interpreters like Barrett,
de Jong, and Dickin, acquire authority from the way that they deliver their
message, and it has long been a truism that messengers speak with `Homeric diction.'
But what is `Homeric diction,' and how might we identify it, especially when
messengers' speeches appear to use the same dialect as other spoken sections
of tragedy.

Part of the difficulty here lies in the limited data available for quantitative
analyses of Athenian drama. As noted above, the entire corpus under analysis here
consists of 44,482 lines, of which we analyzed 18,968 trimeters. Over half of
those lines belong to Euripides' plays; the other half skews slightly Sophoclean.
But this relatively small dataset and its unequal distribution mean that quantitative
analyses are necessarily limited in scope. As Quinn Dombrowski and Patrick Burns
note, languages other than English are at a severe disadvantage when it comes
to computational study.\cite{Dombrowski.Burns2023} Although quantitative study
of drama has been conducted successfully by \cite{Beine2025} and others, those
studies were carried out on larger corpora with a greater concentration of
manual annotations than are currently available for Ancient Greek.

Especially for poetic texts like tragedy, the state of the art lemmatizers
and morphological tools---which have been trained overwhelmingly on prose
treebanks---do not perform sufficiently well for nuanced analysis. This
limitation motivated the present study's focus on meter using a relatively
naïve rules-based parser on a relatively straightforward meter. (By avoiding
anapests and odes, the metrical parser could avoid many of the pitfalls
that would require extensive training and manual annotation/intervention.)
Yet even with these self-imposed guardrails, the present study could not
marshall sufficient data to detect any metrical signature on the part of
character-type speech.

Under-resourced languages must find ways to overcome these barriers if
computational scholars are to be seen on equal footing with our English-
(and Spanish-, French-, German-, Italian-, Chinese-, etc.) language–focused
colleagues. As a start, we need to recognize the labor that goes into annotating
textual sources. The Text Encoding Initiative is something as a barrier here,
as it focuses on physical texts rather than the abstract textual models needed
by editions of drama.\footnote{See \cite{Efer2017}.}

But we must also find ways of posing questions the answers to which are not limited
by the limitations of available data. These methodological problems deserve
scholarly attention, as a richer understanding of under-resourced languages---ancient
or modern---leads to a richer appreciation even of well-resourced fields of inquiry.


\section{Future directions}

Even if we cannot use meter to distinguish messenger-type from other kinds of
speech in tragedy, `Homeric diction' might still be a useful heuristic for
syntax and lexis. To this end, future work will focus on possible lexical
distinctions among different types of speakers in tragedy. Bag-of-words
models offer some reprieve from the limitations of line-oriented studies
(like the present examinations of meter), but successful bag-of-words approaches
depend either on sub-word tokenization or better lemmatization than is currently
available. Studies of syntax are likewise limited by the accuracy of current
morphological tools on poetry.

In a separate vein, future work might focus on network analysis, which would
require additional manual annotations of existing corpora. Such annotations
could be accomplished by a small team of researchers in a reasonably short
amount of time---for instance, by adding `toWhom' attributes to dialogue
tags in TEI-encoded dramas---but could yield interesting results about
internal audiences and attention in the extant plays.

Above all, as this paper has shown, future work should address the
methodological gap between the state-of-the-art models---which focus almost
exclusively on English---and under-resourced langauges like Ancient Greek. By
developing methodologies that can account for the limitations of small
corpora, we can direct efforts to annotate and supplement those corpora while
simultaneously conducting meaningful research.





\section{Declaration on Generative AI}

The author built the metrical parsing library with the help of Anthropic's
Claude Code---and subsequently rewrote most of it without AI assistance. No
AI was used in the calculations or analysis that ultimately make up the bulk
of this paper.


\bibliography{zotero-refs}


\end{document}
